\documentclass[11pt]{article}

\usepackage[utf8]{inputenc}
\usepackage[T1]{fontenc}
\usepackage{geometry}
\usepackage{amsmath}
\usepackage{amssymb}
\usepackage{booktabs}
\usepackage{hyperref}
\usepackage{xurl}
\usepackage{xcolor}
\usepackage{listings}
\usepackage{enumitem}
\usepackage{parskip}
\usepackage{graphicx}
\graphicspath{{../images/}}

\geometry{margin=1in}

\lstset{
  language=Java,
  basicstyle=\ttfamily\small,
  keywordstyle=\color{blue},
  commentstyle=\color{gray},
  stringstyle=\color{teal},
  showstringspaces=false,
  columns=fullflexible,
  frame=single,
  framerule=0.2pt,
  breaklines=true
}

\setlist[itemize]{topsep=0.3em,itemsep=0.2em}
\setlist[enumerate]{topsep=0.3em,itemsep=0.2em}

% Simple per-section source line.
\newcommand{\notesources}[1]{\par\noindent{\small\color{gray}\textit{Sources:} \sloppy #1\par}}

% Unit-wise source strings for this subject (used by slides/notes).
% Path is relative to the lecture `latex/` folder (the compilation CWD).
% OOPS with Java (CSEG1044) — unit-wise sources
%
% These are short, student-facing reference strings intended to be used with:
%   - \slidesources{...}  (in beamer slides)
%   - \notesources{...}   (in lecture notes)
%
% Style inspiration: other subjects in My-Subjects/ use semicolon-separated sources
% and include an "accessed" date for web documentation.

\newcommand{\OOPSUnitISources}{Schildt, \textit{Java: A Beginner's Guide} (10e), McGraw-Hill (Java fundamentals + OOP basics).; Oracle Java Tutorials: Getting Started + Language Basics + Classes and Objects (accessed \today).; Java SE API docs (e.g., \texttt{String}, \texttt{Scanner}) (accessed \today).}

\newcommand{\OOPSUnitIISources}{Schildt, \textit{Java: The Complete Reference} (12e), McGraw-Hill (classes, inheritance, packages, interfaces).; Bloch, \textit{Effective Java} (3e), Addison-Wesley, 2018 (design choices; interfaces vs abstract classes).; Oracle Java Tutorials: Classes and Objects + Interfaces + Packages (accessed \today).; Java SE API docs (e.g., \texttt{Comparable}, \texttt{Runnable}) (accessed \today).}

\newcommand{\OOPSUnitIIISources}{Schildt, \textit{Java: The Complete Reference} (12e) (nested/inner classes, exceptions, threads, I/O).; Oracle Java Tutorials: Nested Classes + Exceptions + Concurrency + Basic I/O (accessed \today).; Java SE API docs (e.g., \texttt{Thread}, \texttt{AutoCloseable}) (accessed \today).}

\newcommand{\OOPSUnitIVSources}{Schildt, \textit{Java: The Complete Reference} (12e) (generics, lambdas, Swing, JDBC).; Bloch, \textit{Effective Java} (3e) (generics + lambdas).; Oracle Java Tutorials: Generics + Lambda Expressions + Swing + JDBC Basics (accessed \today).; Java SE API docs (e.g., \texttt{java.util.function}, \texttt{javax.swing}, \texttt{java.sql}) (accessed \today).}

\newcommand{\OOPSUnitVSources}{Schildt, \textit{Java: The Complete Reference} (12e) (Collections Framework + wrapper classes).; Oracle Java docs: Collections Framework overview (accessed \today).; Java SE API docs (\texttt{java.util} collections; wrapper classes in \texttt{java.lang}) (accessed \today).}

\newcommand{\OOPSUnitVISources}{Git SCM: \textit{Pro Git} (2e) (version control workflow), accessed \today.; GitHub Docs (repositories, issues, pull requests), accessed \today.; Oracle Java Tutorials: Swing + JDBC (accessed \today).; JUnit 5 User Guide (accessed \today).; Apache Maven Project: Getting Started (accessed \today).; Gradle User Manual: Getting Started (accessed \today).; UML class diagrams: UML 2.x overview/spec (accessed \today).}

\newcommand{\OOPSMiscWhyInterfacesSources}{Schildt, \textit{Java: The Complete Reference} (12e) (interfaces).; Bloch, \textit{Effective Java} (3e) (prefer interfaces where appropriate).; Oracle Java Tutorials: Interfaces (accessed \today).; Java SE API docs (e.g., \texttt{Comparable}, \texttt{Runnable}) (accessed \today).}



\title{Object Oriented Programming (CSEG1044)\\Unit 06 -- Lecture 06 Notes\\Testing + Build/Artifacts + Final Demo}
\author{Tofik Ali}
\date{\today}

\begin{document}
\maketitle
\tableofcontents

\section{Lecture Overview}
This is the final capstone lecture.
Your goal is not to add new features endlessly --- your goal is to make the project \textbf{demo-ready}:
\begin{itemize}
  \item core flows work reliably,
  \item mistakes are handled gracefully (validation + errors),
  \item the project builds into a runnable artifact,
  \item you have a clear demo script and evidence of teamwork.
\end{itemize}
\notesources{\OOPSUnitVISources}

\section{Syllabus Alignment (Unit 06)}
\textbf{Unit 06:} testing (JUnit), build tool + packaging, final demo/viva

\section{Learning Outcomes}
After this lecture, you should be able to:
\begin{itemize}
  \item Write unit tests for your service/business rules using JUnit 5.
  \item Perform an integration test checklist for DAO + DB.
  \item Package the project as a runnable artifact (Jar) using Maven/Gradle (or a simple build script).
  \item Complete Deliverable-5: final demo/viva + release artifact + short report.
\end{itemize}

\section{Key Terms (Quick)}
\begin{itemize}
  \item \textbf{Unit test}: tests one class/function in isolation (fast, deterministic).
  \item \textbf{Integration test}: tests multiple parts together (DAO + DB, UI + service).
  \item \textbf{Test pyramid}: many unit tests, fewer integration tests, minimal UI/manual tests.
  \item \textbf{Assertion}: check that expected result is correct (\texttt{assertEquals}, \texttt{assertThrows}).
  \item \textbf{Fixture}: setup data/objects for tests (\texttt{@BeforeEach}).
  \item \textbf{Artifact}: build output you can run/distribute (Jar).
  \item \textbf{Release}: versioned snapshot of your project (tagged commit + artifact).
\end{itemize}

\section{Detailed Notes}
\subsection{What should you test? (capstone strategy)}
\textbf{Best unit-test targets:}
\begin{itemize}
  \item service-layer validation and business rules (``no overlap'', ``capacity $>$ 0'')
  \item small pure functions (parsers, formatters)
\end{itemize}

\textbf{Harder to unit test (do minimal/manual):}
\begin{itemize}
  \item Swing UI (usually tested manually in this course)
  \item DAO mapping (better as integration tests with a test database)
\end{itemize}

\textbf{Why?} Service rules are where bugs hurt the most, and they are easy to test quickly without a GUI or DB.

\subsection{JUnit 5 basics (how a unit test looks)}
JUnit tests have a simple structure:
\begin{itemize}
  \item Arrange: prepare objects
  \item Act: call the method
  \item Assert: check result/exception
\end{itemize}

\begin{lstlisting}
import static org.junit.jupiter.api.Assertions.*;
import org.junit.jupiter.api.Test;

class SimpleTest {
    @Test
    void endMustBeAfterStart() {
        ValidationException ex = assertThrows(
            ValidationException.class,
            () -> BookingRules.validateTime(10, 10)
        );
        assertTrue(ex.getMessage().contains("after"));
    }
}
\end{lstlisting}

\textbf{Tip:} test names should describe behavior: \texttt{rejectsEndBeforeStart}.

\subsection{Design for testability: depend on interfaces}
If your service depends directly on a JDBC DAO, tests become slow and fragile.
Instead:
\begin{itemize}
  \item define an interface (contract), e.g., \texttt{BookingRepository},
  \item implement JDBC repository for real app,
  \item implement an in-memory fake for unit tests.
\end{itemize}

\subsection{Integration test checklist (DAO + DB)}
Use a separate test DB/schema (never your real production data).
\begin{enumerate}
  \item Run \texttt{schema.sql} on a clean database.
  \item Run DAO insert for one entity.
  \item Run DAO list/find and verify data is correct.
  \item Run update and verify changes.
  \item Run delete and verify row is gone.
  \item Test one ``bad input'' case (FK constraint violation, duplicate UNIQUE).
\end{enumerate}

\subsection{Build tool + packaging (Maven/Gradle idea)}
\textbf{Why use a build tool?}
\begin{itemize}
  \item consistent compilation and testing for everyone in team,
  \item dependency management (JDBC driver, JUnit),
  \item easy packaging into a Jar.
\end{itemize}

\textbf{Maven (common commands):}
\begin{itemize}
  \item \texttt{mvn test} --- run unit tests
  \item \texttt{mvn package} --- build Jar in \texttt{target/}
\end{itemize}

\textbf{Gradle (common commands):}
\begin{itemize}
  \item \texttt{gradle test}
  \item \texttt{gradle build}
\end{itemize}

\textbf{Packaging note:} if your Jar needs external libraries (like MySQL connector),
you either run with a classpath, or build a ``fat Jar'' that bundles dependencies.

\notesources{\OOPSUnitVISources}

\section{Worked Example: unit-test a BookingService without DB}
This example shows the key pattern for capstone testing: service depends on an interface, not JDBC.

\subsection{Repository interface + in-memory fake}
\begin{lstlisting}
import java.util.*;

interface BookingRepository {
    boolean hasOverlap(int resourceId, long start, long end);
    void save(Booking b);
    int count();
}

class InMemoryBookingRepository implements BookingRepository {
    private final List<Booking> data = new ArrayList<>();

    public boolean hasOverlap(int resourceId, long start, long end) {
        for (Booking b : data) {
            if (b.getResourceId() != resourceId) continue;
            boolean overlaps = start < b.getEnd() && end > b.getStart();
            if (overlaps) return true;
        }
        return false;
    }

    public void save(Booking b) { data.add(b); }
    public int count() { return data.size(); }
}
\end{lstlisting}

\subsection{Service (rules) + tests}
\begin{lstlisting}
class BookingService {
    private final BookingRepository repo;
    BookingService(BookingRepository repo) { this.repo = repo; }

    public void createBooking(int resourceId, String user, long start, long end) {
        if (user == null || user.trim().isEmpty())
            throw new ValidationException("User is required.");
        if (end <= start)
            throw new ValidationException("End must be after start.");
        if (repo.hasOverlap(resourceId, start, end))
            throw new ValidationException("Booking overlaps existing booking.");
        repo.save(new Booking(resourceId, user.trim(), start, end));
    }
}
\end{lstlisting}

\begin{lstlisting}
import static org.junit.jupiter.api.Assertions.*;
import org.junit.jupiter.api.*;

class BookingServiceTest {
    private InMemoryBookingRepository repo;
    private BookingService service;

    @BeforeEach
    void setup() {
        repo = new InMemoryBookingRepository();
        service = new BookingService(repo);
    }

    @Test
    void rejectsEndBeforeStart() {
        assertThrows(ValidationException.class,
            () -> service.createBooking(1, "A", 10, 10));
    }

    @Test
    void rejectsOverlap() {
        service.createBooking(1, "A", 10, 20);
        assertThrows(ValidationException.class,
            () -> service.createBooking(1, "B", 15, 25));
    }

    @Test
    void acceptsNonOverlap() {
        assertDoesNotThrow(() -> service.createBooking(1, "A", 10, 20));
        assertDoesNotThrow(() -> service.createBooking(1, "B", 21, 30));
        assertEquals(2, repo.count());
    }
}
\end{lstlisting}

\section{Deliverable-5 checklist (final submission)}
\begin{itemize}
  \item Runnable artifact: Jar (plus any config instructions).
  \item \texttt{schema.sql} + DB setup instructions.
  \item Basic unit tests for service rules (JUnit) + evidence they pass.
  \item README: how to run, demo steps, team members + responsibilities.
  \item Short report (1--2 pages): design summary + what was implemented.
\end{itemize}

\section{Final Demo Script (suggested)}
\begin{enumerate}
  \item Start DB and ensure schema exists.
  \item Run the application.
  \item Show one end-to-end use case: Add Resource $\rightarrow$ verify in list.
  \item Show booking flow: create booking $\rightarrow$ list bookings.
  \item Show validation: attempt an invalid booking (end before start) and show message.
  \item Briefly show code structure: UI/service/DAO packages.
  \item Show tests running (optional but impressive): \texttt{mvn test}.
\end{enumerate}

\section{Checkpoint Questions (with sample answers)}
\textbf{Q1.} Why unit test service logic instead of Swing UI?\par
\textbf{Sample answer:} Service logic is easy to test quickly and deterministically; UI tests are slower and fragile.

\vspace{0.6em}
\textbf{Q2.} What does \texttt{assertThrows} test?\par
\textbf{Sample answer:} It verifies that a method call throws an expected exception for an invalid case.

\vspace{0.6em}
\textbf{Q3.} Why use a build tool?\par
\textbf{Sample answer:} It standardizes builds/tests, manages dependencies, and produces a runnable artifact reliably for the team.

\end{document}
